% Source: https://artofproblemsolving.com/community/c3h1745554_system_of_nonlinear_equations
% Stable identity: aops-topic-1745554-q01
Let $a=x^3$, $b=y^3$; since cube roots of reals are unique, real pairs $(x,y)$ correspond exactly to real pairs $(a,b)$ with $a^2+b^2=6969$, $a^3+b^3=69$. Put $s=a+b$, $p=ab$. Then $p=\dfrac{s^2-6969}{2}$ and
\[
69=a^3+b^3=s^3-3ps=s^3-\dfrac{3s(s^2-6969)}{2}\iff f(s):=s^3-20907s+138=0 .
\]
For $a,b$ to be real we need $s^2-4p\ge0$, i.e.\ $s^2\le13938$. Let $m=\sqrt{13938}$. We have $f(-m)=6969m+138>0$ and $f(m)=-6969m+138<0$. The derivative $3s^2-20907$ vanishes at $\pm\sqrt{6969}$, both inside $(-m,m)$; $f$ increases on $[-m,-\sqrt{6969}]$ (so it stays positive there), decreases on $[-\sqrt{6969},\sqrt{6969}]$, and increases on $[\sqrt{6969},m]$ up to the negative value $f(m)$. Hence $f$ has exactly one root $s_0$ with $s_0^2\le13938$ (in fact $s_0\approx0.0066$), and then $s_0^2<13938$, so $a\ne b$ are the two distinct real roots of $t^2-s_0t+p_0$.

The real solutions are therefore $(a,b)$ and $(b,a)$, giving the ratios $\dfrac {x}{y}=\sqrt[3]{\dfrac {a}{b}}$ and $\sqrt[3]{\dfrac {b}{a}}$. Neither $a$ nor $b$ is $0$ (that would force $6969^{3}=69^2$), and these two ratios are reciprocals, equal only if $\dfrac {a}{b}=\pm1$; $a=b$ is excluded and $a=-b$ would give $a^3+b^3=0\ne69$. So $\dfrac {x}{y}$ takes exactly $\boxed{2}$ values.
